\documentclass[11pt,a4paper]{article}

% Packages
\usepackage[utf8]{inputenc}
\usepackage[T1]{fontenc}
\usepackage{amsmath,amssymb}
\usepackage{graphicx}
\usepackage{xcolor}
\usepackage{listings}
\usepackage{hyperref}
\usepackage[margin=1in]{geometry}
\usepackage{enumitem}
\usepackage{tocloft}
\usepackage{booktabs}
\usepackage{fancyhdr}
\usepackage{float}

% Colors
\definecolor{codegreen}{rgb}{0,0.6,0}
\definecolor{codegray}{rgb}{0.5,0.5,0.5}
\definecolor{codepurple}{rgb}{0.58,0,0.82}
\definecolor{backcolour}{rgb}{0.95,0.95,0.92}
\definecolor{darkblue}{rgb}{0.0,0.0,0.6}
\definecolor{darkgreen}{rgb}{0.0,0.4,0.0}

% Listing configuration for C++ code
\lstdefinestyle{mystyle}{
    language=C++,
    backgroundcolor=\color{backcolour},   
    commentstyle=\color{codegreen},
    keywordstyle=\color{darkblue}\bfseries,
    numberstyle=\tiny\color{codegray},
    stringstyle=\color{codepurple},
    basicstyle=\ttfamily\footnotesize,
    breakatwhitespace=false,         
    breaklines=true,                 
    captionpos=b,                    
    keepspaces=true,                 
    numbers=left,                    
    numbersep=5pt,                  
    showspaces=false,                
    showstringspaces=false,
    showtabs=false,                  
    tabsize=2,
    frame=single,
    morekeywords={vector,cv,Mat,KeyPoint,MapPoint,Sophus,SE3,SE3f,Eigen,Matrix3f,Vector3f,Frame,Tracking,System,Atlas,KeyFrame}
}
\lstset{style=mystyle}

% Hyperref setup
\hypersetup{
    colorlinks=true,
    linkcolor=darkblue,
    filecolor=magenta,      
    urlcolor=cyan,
    citecolor=darkgreen,
}

% Header/Footer
\pagestyle{fancy}
\fancyhf{}
\rhead{ORB-SLAM3 Code Story}
\lhead{Mohamed Alimoussa}
\cfoot{\thepage}

% Title
\title{\textbf{The Story of ORB-SLAM3}\\
\large{A Complete Code Walkthrough from Main to MapPoint}\\[1em]
\normalsize{Understanding How ORB-SLAM3 Processes Images and Builds Maps}}
\author{Mohamed Alimoussa}
\date{January 2026}

\begin{document}

\maketitle

\begin{abstract}
This document tells the complete story of how ORB-SLAM3 works, following the code execution path from the moment you run the program to the creation of 3D maps. We trace the journey of an image as it flows through the system: from being loaded at the command line, through feature extraction, pose estimation, map building, and loop closure. This narrative approach helps you understand not just \textit{what} each component does, but \textit{why} it exists and \textit{how} it connects to everything else.
\end{abstract}

\tableofcontents
\newpage

% Include chapter files
%==============================================================================
\section{Prologue: The Cast of Characters}
%==============================================================================

Before we begin our story, let's meet the main characters. ORB-SLAM3 is built around a handful of key classes that work together like actors on a stage.

\subsection{The Main Characters}

\begin{table}[H]
\centering
\begin{tabular}{@{}lll@{}}
\toprule
\textbf{Class} & \textbf{Role} & \textbf{File} \\ 
\midrule
\texttt{System} & The Director & \texttt{System.h/.cc} \\
\texttt{Tracking} & The Lead Actor & \texttt{Tracking.h/.cc} \\
\texttt{LocalMapping} & The Builder & \texttt{LocalMapping.h/.cc} \\
\texttt{LoopClosing} & The Proofreader & \texttt{LoopClosing.h/.cc} \\
\texttt{Atlas} & The Library & \texttt{Atlas.h/.cc} \\
\texttt{Frame} & The Photo & \texttt{Frame.h/.cc} \\
\texttt{KeyFrame} & The Important Photo & \texttt{KeyFrame.h/.cc} \\
\texttt{MapPoint} & The 3D Landmark & \texttt{MapPoint.h/.cc} \\
\bottomrule
\end{tabular}
\caption{The main classes in ORB-SLAM3 and their roles}
\end{table}

\subsection{The Supporting Cast}

\begin{itemize}[leftmargin=*]
    \item \textbf{ORBextractor}: Finds interesting points in images (or \texttt{ONNXFeatureExtractor} for XFeat)
    \item \textbf{ORBVocabulary}: A dictionary for recognizing places (used for loop closure)
    \item \textbf{GeometricCamera}: Knows how to project 3D points onto the image
    \item \textbf{g2o}: The optimization library that refines everything
    \item \textbf{Viewer}: Draws the map so you can see what's happening
\end{itemize}

\subsection{How They Work Together}

Think of ORB-SLAM3 like a film production:

\begin{enumerate}
    \item \textbf{The Director (System)} doesn't act, but coordinates everyone
    \item \textbf{The Lead Actor (Tracking)} is always on stage, processing every image
    \item \textbf{The Builder (LocalMapping)} works backstage, building the set
    \item \textbf{The Proofreader (LoopClosing)} checks for mistakes and fixes them
    \item \textbf{The Photos (Frame/KeyFrame)} are the snapshots that document everything
    \item \textbf{The Landmarks (MapPoint)} are the 3D decorations that make the scene real
\end{enumerate}

%==============================================================================
\section{Act I: The Beginning (\texttt{mono\_euroc.cc})}
%==============================================================================

Our story begins in the \texttt{Examples/Monocular/} folder. This is the \textbf{entry point}---the \texttt{main()} function where everything starts.

\subsection{Scene 1: The Command Line}

When you run ORB-SLAM3 from the terminal:
\begin{lstlisting}[language=bash,numbers=none]
./mono_euroc ../Vocabulary/ORBvoc.txt TUM_mono.yaml /path/to/MH_01 timestamps.txt
\end{lstlisting}

The program expects these arguments:
\begin{enumerate}
    \item \textbf{Vocabulary file}: The ORB dictionary for place recognition
    \item \textbf{Settings file}: Camera calibration, feature parameters, etc.
    \item \textbf{Image folder}: Where the images are stored
    \item \textbf{Timestamps file}: When each image was captured
\end{enumerate}

\subsection{Scene 2: Loading the Images}

The first thing the program does is load the list of images:

\begin{lstlisting}[caption={Loading images in mono\_euroc.cc (simplified)},label=lst:loadimages]
// Load all image paths and timestamps
vector<string> vstrImageFilenames;
vector<double> vTimestamps;

LoadImages("/path/to/MH_01/mav0/cam0/data", 
           "timestamps.txt",
           vstrImageFilenames, 
           vTimestamps);

// Now we have:
// vstrImageFilenames[0] = ".../1403636579763555584.png"
// vstrImageFilenames[1] = ".../1403636579813555456.png"
// ...
\end{lstlisting}

\textbf{Key insight:} At this point, we haven't done any SLAM yet. We're just preparing a list of ``what images to process''.

\subsection{Scene 3: Creating the SLAM System}

This is the most important line in the entire file:

\begin{lstlisting}[caption={Creating the SLAM System}]
ORB_SLAM3::System SLAM(argv[1], argv[2], ORB_SLAM3::System::MONOCULAR, true);
//                     ^vocab    ^config  ^sensor type              ^use viewer
\end{lstlisting}

When this line executes, \textbf{the entire SLAM machinery is created}:
\begin{itemize}
    \item The vocabulary is loaded from disk (takes a few seconds)
    \item Camera calibration is parsed from the YAML file
    \item Three threads are spawned:
    \begin{enumerate}
        \item \textbf{Tracking} (runs in the main thread)
        \item \textbf{LocalMapping} (background thread)
        \item \textbf{LoopClosing} (background thread)
    \end{enumerate}
    \item The Viewer window opens (if enabled)
\end{itemize}

\textit{What happens inside the \texttt{System} constructor? We'll explore this in the next section.}

\subsection{Scene 4: The Main Loop}

After initialization, we enter the main loop---the heartbeat of SLAM:

\begin{lstlisting}[caption={The main processing loop (simplified)}]
for(int ni = 0; ni < nImages; ni++)
{
    // 1. Read the image from disk
    cv::Mat im = cv::imread(vstrImageFilenames[ni], cv::IMREAD_UNCHANGED);
    double timestamp = vTimestamps[ni];
    
    // 2. THE MAGIC LINE: Feed image to SLAM
    SLAM.TrackMonocular(im, timestamp);
    
    // 3. (Optional) Wait to simulate real-time
    // usleep(...);
}
\end{lstlisting}

\textbf{Key insight:} The entire complexity of SLAM is hidden behind one function call: \texttt{TrackMonocular()}. This function takes a raw image and returns the camera pose. Everything else---feature extraction, matching, optimization, map building---happens inside.

\subsection{Scene 5: Cleanup}

After all images are processed:
\begin{lstlisting}[caption={Shutting down}]
SLAM.Shutdown();  // Stop all threads gracefully

// Save trajectory for evaluation
SLAM.SaveTrajectoryEuRoC("CameraTrajectory.txt");
SLAM.SaveKeyFrameTrajectoryEuRoC("KeyFrameTrajectory.txt");
\end{lstlisting}

\subsection{Summary: The Entry Point Flow}

\begin{verbatim}
main()
  |
  +-- Parse command line arguments
  |
  +-- LoadImages() --> Fill vstrImageFilenames, vTimestamps
  |
  +-- Create System SLAM(...) --> Spawns 3 threads
  |
  +-- for each image:
  |       SLAM.TrackMonocular(im, timestamp)
  |
  +-- SLAM.Shutdown()
  |
  +-- SLAM.SaveTrajectory...()
\end{verbatim}

%==============================================================================
\section{Act II: The Brain (\texttt{System.h/.cc})}
%==============================================================================

The \texttt{System} class is the \textbf{orchestrator}. It doesn't do any computation itself, but it creates and coordinates all the components.

\subsection{What Does the System Own?}

Looking at \texttt{System.h}, we see the System holds pointers to everything:

\begin{lstlisting}[caption={System's private members (simplified)}]
private:
    eSensor mSensor;              // MONOCULAR, STEREO, RGBD, etc.
    
    // The knowledge base
    ORBVocabulary* mpVocabulary;  // Visual vocabulary for recognition
    KeyFrameDatabase* mpKeyFrameDatabase;  // For loop detection
    Atlas* mpAtlas;               // Collection of all maps
    
    // The workers
    Tracking* mpTracker;          // Estimates camera pose
    LocalMapping* mpLocalMapper;  // Builds the map
    LoopClosing* mpLoopCloser;    // Fixes drift
    
    // Visualization
    Viewer* mpViewer;
    FrameDrawer* mpFrameDrawer;
    MapDrawer* mpMapDrawer;
    
    // Threads
    std::thread* mptLocalMapping;
    std::thread* mptLoopClosing;
    std::thread* mptViewer;
\end{lstlisting}

\subsection{The System Constructor: Birth of SLAM}

When you call \texttt{System(vocabPath, settingsPath, MONOCULAR, true)}, this happens:

\begin{enumerate}
    \item \textbf{Load Vocabulary}: Read the ORB dictionary from disk (~100MB, takes a few seconds)
    \item \textbf{Create Atlas}: The ``library'' that will store all maps and map points
    \item \textbf{Create KeyFrameDatabase}: For finding similar places (loop closure)
    \item \textbf{Create Drawers}: For visualization
    \item \textbf{Create Tracking}: The main worker (we'll explore this in the next section)
    \item \textbf{Create LocalMapping}: Spawn a background thread
    \item \textbf{Create LoopClosing}: Spawn another background thread
    \item \textbf{Create Viewer}: Spawn the visualization thread
\end{enumerate}

\subsection{The Sensor Types}

The System supports multiple sensor configurations:

\begin{lstlisting}[caption={Sensor types enum}]
enum eSensor {
    MONOCULAR = 0,     // Single camera
    STEREO = 1,        // Two cameras (left/right)
    RGBD = 2,          // RGB + Depth camera (Kinect, RealSense)
    IMU_MONOCULAR = 3, // Mono + IMU
    IMU_STEREO = 4,    // Stereo + IMU
    IMU_RGBD = 5       // RGBD + IMU
};
\end{lstlisting}

\subsection{The Public Interface}

The System exposes a simple interface to the outside world:

\begin{lstlisting}[caption={System's main interface}]
// Feed images to the system (returns camera pose)
Sophus::SE3f TrackMonocular(const cv::Mat &im, const double &timestamp);
Sophus::SE3f TrackStereo(const cv::Mat &imLeft, const cv::Mat &imRight, ...);
Sophus::SE3f TrackRGBD(const cv::Mat &im, const cv::Mat &depthmap, ...);

// Control
void ActivateLocalizationMode();  // Stop mapping, only track
void DeactivateLocalizationMode();
void Reset();
void Shutdown();

// Output
void SaveTrajectoryTUM(const string &filename);
void SaveTrajectoryEuRoC(const string &filename);
\end{lstlisting}

\textbf{Key insight:} The System class follows the \textit{Facade} design pattern. It hides all the complexity behind a simple interface. You don't need to know about Frames, KeyFrames, or MapPoints to use it.

\subsection{What Happens in \texttt{TrackMonocular}?}

When you call \texttt{SLAM.TrackMonocular(im, timestamp)}:

\begin{lstlisting}[caption={TrackMonocular simplified flow}]
Sophus::SE3f System::TrackMonocular(const cv::Mat &im, const double &timestamp)
{
    // 1. Convert to grayscale if needed
    cv::Mat imGray;
    if(im.channels() == 3)
        cv::cvtColor(im, imGray, cv::COLOR_BGR2GRAY);
    else
        imGray = im;
    
    // 2. Delegate to Tracking (the real worker!)
    Sophus::SE3f Tcw = mpTracker->GrabImageMonocular(imGray, timestamp, "");
    
    // 3. Update state for external queries
    mTrackingState = mpTracker->mState;
    mTrackedMapPoints = mpTracker->mCurrentFrame.mvpMapPoints;
    
    return Tcw;  // Return the camera pose
}
\end{lstlisting}

So the System is essentially a thin wrapper. The real work happens in \texttt{Tracking}.

\subsection{The Multi-Threading Model}

ORB-SLAM3 uses a producer-consumer pattern:

\begin{itemize}
    \item \textbf{Tracking (Main Thread)}: Produces KeyFrames, consumes images
    \item \textbf{LocalMapping (Thread 1)}: Consumes KeyFrames, produces MapPoints
    \item \textbf{LoopClosing (Thread 2)}: Consumes KeyFrames, produces corrections
    \item \textbf{Viewer (Thread 3)}: Consumes map state, produces visualization
\end{itemize}

Communication happens through:
\begin{itemize}
    \item Shared pointers to KeyFrames and MapPoints
    \item Thread-safe queues
    \item Mutex-protected data access
\end{itemize}

%==============================================================================
\section{Act III: The Hero (\texttt{Tracking.h/.cc})}
%==============================================================================

The \texttt{Tracking} class is where the magic happens. It takes raw images and produces camera poses.

\subsection{The Tracking States}

Tracking operates like a state machine:

\begin{lstlisting}[caption={Tracking states}]
enum eTrackingState {
    SYSTEM_NOT_READY = -1,  // Still initializing
    NO_IMAGES_YET = 0,      // Haven't received any images
    NOT_INITIALIZED = 1,    // Trying to initialize the map
    OK = 2,                 // Normal operation
    RECENTLY_LOST = 3,      // Just lost tracking (trying to recover)
    LOST = 4,               // Completely lost
    OK_KLT = 5              // Tracking with KLT (fallback)
};
\end{lstlisting}

\textbf{The lifecycle:}
\begin{enumerate}
    \item Start in \texttt{NO\_IMAGES\_YET}
    \item First image $\rightarrow$ transition to \texttt{NOT\_INITIALIZED}
    \item After successful initialization $\rightarrow$ transition to \texttt{OK}
    \item If tracking fails $\rightarrow$ \texttt{RECENTLY\_LOST} $\rightarrow$ try relocalization
    \item If relocalization fails $\rightarrow$ \texttt{LOST}
\end{enumerate}

\subsection{The Main Entry: GrabImageMonocular}

\texttt{GrabImageMonocular} is called by System for each new image:

\begin{lstlisting}[caption={GrabImageMonocular simplified}]
Sophus::SE3f Tracking::GrabImageMonocular(const cv::Mat &im, 
                                          const double &timestamp, 
                                          string filename)
{
    // 1. Store the image
    mImGray = im;
    
    // 2. Create a Frame (feature extraction happens here!)
    mCurrentFrame = Frame(mImGray, timestamp, mpORBextractorLeft, 
                          mpORBVocabulary, mpCamera, mDistCoef, 
                          mbf, mThDepth, &mLastFrame, mImuCalib);
    
    // 3. The main tracking logic
    Track();
    
    // 4. Return the estimated pose
    return mCurrentFrame.GetPose();
}
\end{lstlisting}

\subsection{Inside \texttt{Track()}: The Heart of SLAM}

The \texttt{Track()} function is a 500+ line behemoth, but its structure is simple:

\begin{lstlisting}[caption={Track() simplified structure}]
void Tracking::Track()
{
    // PHASE 1: INITIALIZATION
    if(mState == NOT_INITIALIZED)
    {
        MonocularInitialization();  // Try to initialize the map
        return;
    }
    
    // PHASE 2: POSE ESTIMATION (if initialized)
    bool bOK = false;
    
    // First attempt: Use motion model
    if(mbVelocity)
        bOK = TrackWithMotionModel();
    
    // Fallback: Match against reference keyframe
    if(!bOK)
        bOK = TrackReferenceKeyFrame();
    
    // If still failing: Try relocalization
    if(!bOK)
        bOK = Relocalization();
    
    // PHASE 3: TRACK LOCAL MAP (refine pose using more points)
    if(bOK)
        bOK = TrackLocalMap();
    
    // PHASE 4: KEYFRAME DECISION
    if(bOK && NeedNewKeyFrame())
        CreateNewKeyFrame();
    
    // Update state
    mState = bOK ? OK : LOST;
}
\end{lstlisting}

\subsection{Phase 1: Monocular Initialization}

Monocular SLAM has a chicken-and-egg problem:
\begin{itemize}
    \item To estimate the camera pose, we need 3D points
    \item To triangulate 3D points, we need two camera poses
\end{itemize}

The solution: \textbf{Two-View Reconstruction}

\begin{lstlisting}[caption={MonocularInitialization simplified}]
void Tracking::MonocularInitialization()
{
    // Step 1: First frame becomes the "initial frame"
    if(mState == NO_IMAGES_YET)
    {
        mInitialFrame = Frame(mCurrentFrame);
        return;
    }
    
    // Step 2: Match features between initial and current frame
    vector<int> vMatches;
    int nMatches = ORBmatcher::SearchForInitialization(
        mInitialFrame, mCurrentFrame, vMatches);
    
    if(nMatches < 100)
        return;  // Not enough matches, wait for better frame
    
    // Step 3: Compute Homography and Fundamental matrix
    // (two competing geometric models)
    cv::Mat H, F;
    float SH, SF;  // Scores for each model
    ComputeH21(vMatches, SH, H);
    ComputeF21(vMatches, SF, F);
    
    // Step 4: Choose the better model and recover pose
    cv::Mat R, t;
    vector<cv::Point3f> vP3D;  // Triangulated 3D points
    
    float RH = SH / (SH + SF);  // "Homographyness"
    
    if(RH > 0.4)
        ReconstructH(H, R, t, vP3D);  // Planar scene
    else
        ReconstructF(F, R, t, vP3D);  // General scene
    
    // Step 5: Create the initial map
    CreateInitialMapMonocular();
    
    mState = OK;  // Initialization complete!
}
\end{lstlisting}

\textbf{Key insight:} After initialization, we have:
\begin{itemize}
    \item Two KeyFrames (the initial frame and current frame)
    \item A set of 3D MapPoints (triangulated from the two views)
    \item An arbitrary scale (monocular SLAM cannot recover absolute scale)
\end{itemize}

\subsection{Phase 2: Pose Estimation}

Once initialized, we estimate the camera pose for each new frame.

\subsubsection{TrackWithMotionModel}

This is the fastest method. We assume the camera is moving at constant velocity:

\begin{lstlisting}[caption={TrackWithMotionModel concept}]
bool Tracking::TrackWithMotionModel()
{
    // 1. Predict pose using velocity
    Sophus::SE3f Tcw_predicted = mVelocity * mLastFrame.GetPose();
    mCurrentFrame.SetPose(Tcw_predicted);
    
    // 2. Project MapPoints from last frame into current frame
    // 3. Search for matches near projected locations
    int nMatches = ORBmatcher::SearchByProjection(mCurrentFrame, mLastFrame);
    
    if(nMatches < 20)
        return false;
    
    // 4. Optimize pose using matched points
    Optimizer::PoseOptimization(&mCurrentFrame);
    
    return true;
}
\end{lstlisting}

\subsubsection{TrackReferenceKeyFrame}

Fallback when motion model fails:

\begin{lstlisting}[caption={TrackReferenceKeyFrame concept}]
bool Tracking::TrackReferenceKeyFrame()
{
    // 1. Compute Bag-of-Words for current frame
    mCurrentFrame.ComputeBoW();
    
    // 2. Match against reference keyframe using BoW
    int nMatches = ORBmatcher::SearchByBoW(mpReferenceKF, mCurrentFrame);
    
    if(nMatches < 15)
        return false;
    
    // 3. Use reference keyframe pose as initial guess
    mCurrentFrame.SetPose(mpReferenceKF->GetPose());
    
    // 4. Optimize pose
    Optimizer::PoseOptimization(&mCurrentFrame);
    
    return true;
}
\end{lstlisting}

\subsection{Phase 3: Track Local Map}

After initial pose estimation, we refine it using more MapPoints:

\begin{lstlisting}[caption={TrackLocalMap concept}]
bool Tracking::TrackLocalMap()
{
    // 1. Find all keyframes that share MapPoints with current frame
    UpdateLocalKeyFrames();
    
    // 2. Collect all MapPoints visible from those keyframes
    UpdateLocalPoints();  // mvpLocalMapPoints
    
    // 3. Project local MapPoints into current frame
    //    Search for matches
    SearchLocalPoints();
    
    // 4. Final pose optimization with all matches
    Optimizer::PoseOptimization(&mCurrentFrame);
    
    // 5. Count inliers
    mnMatchesInliers = CountInliers();
    
    return (mnMatchesInliers > 30);
}
\end{lstlisting}

\textbf{Key insight:} Local Map Tracking uses many more points than the initial pose estimation, giving a more accurate and robust pose.

\subsection{Phase 4: KeyFrame Decision}

Not every frame becomes a KeyFrame. We only insert new KeyFrames when:

\begin{lstlisting}[caption={NeedNewKeyFrame simplified}]
bool Tracking::NeedNewKeyFrame()
{
    // 1. Minimum time since last keyframe
    if(mCurrentFrame.mnId < mnLastKeyFrameId + mMinFrames)
        return false;
    
    // 2. Local mapping must not be busy
    if(mpLocalMapper->isStopped() || mpLocalMapper->stopRequested())
        return false;
    
    // 3. Significant change in observation
    float refRatio = 0.9f;  // 90% of reference KF points still visible
    int nTrackedClose = CountClosePoints();
    
    bool c1 = (mnMatchesInliers < mpReferenceKF->TrackedMapPoints * refRatio);
    bool c2 = (nTrackedClose < 100);  // Few close points
    
    return (c1 || c2);
}
\end{lstlisting}

\textbf{Key insight:} KeyFrames are inserted sparingly to avoid redundancy, but frequently enough to maintain map quality.

\subsection{The Key Member Variables}

\begin{lstlisting}[caption={Important Tracking members}]
// The two most important frames
Frame mCurrentFrame;   // The frame being processed NOW
Frame mLastFrame;      // The previous frame

// Tracking state
eTrackingState mState;

// Current matches
int mnMatchesInliers;

// Local map
KeyFrame* mpReferenceKF;
vector<KeyFrame*> mvpLocalKeyFrames;
vector<MapPoint*> mvpLocalMapPoints;

// Motion model
Sophus::SE3f mVelocity;
bool mbVelocity;

// Feature extraction
ORBextractor* mpORBextractorLeft;
ONNXFeatureExtractor* mpXFeatExtractor;  // Your XFeat mode
\end{lstlisting}

%==============================================================================
\section{Act IV: The Photo (\texttt{Frame.h/.cc})}
%==============================================================================

The \texttt{Frame} class represents a single image capture. It is the fundamental data unit that flows through the system.

\subsection{What Does a Frame Contain?}

\begin{lstlisting}[caption={Key Frame members}]
class Frame {
    // Identity
    long unsigned int mnId;       // Unique frame ID
    double mTimeStamp;            // When captured
    
    // Features (the interesting points)
    int N;                        // Number of keypoints
    vector<cv::KeyPoint> mvKeys;  // 2D keypoint locations
    vector<cv::KeyPoint> mvKeysUn;  // Undistorted keypoints
    cv::Mat mDescriptors;         // Feature descriptors (Nx32 for ORB)
    
    // 3D associations
    vector<MapPoint*> mvpMapPoints;  // Pointer to 3D point for each keypoint
    vector<bool> mvbOutlier;      // Outlier flags
    
    // Camera
    cv::Mat mK;                   // Calibration matrix
    cv::Mat mDistCoef;            // Distortion coefficients
    GeometricCamera* mpCamera;    // Camera model
    
    // Pose
    Sophus::SE3f mTcw;            // Camera pose (world to camera)
    bool mbHasPose;               // Is pose valid?
    
    // Optimization structures
    vector<size_t> mGrid[64][48]; // Spatial grid for fast search
    DBoW2::BowVector mBowVec;     // Bag-of-Words representation
};
\end{lstlisting}

\subsection{The Frame Constructor Journey}

When a new Frame is created, this happens:

\begin{enumerate}
    \item \textbf{Assign ID}: \texttt{mnId = nNextId++}
    \item \textbf{Extract Features}: Call ORBextractor (or XFeat)
    \item \textbf{Undistort Keypoints}: Correct for lens distortion
    \item \textbf{Initialize Associations}: \texttt{mvpMapPoints} all set to NULL
    \item \textbf{Assign to Grid}: Organize keypoints into spatial cells
    \item \textbf{Setup Scale Pyramid}: Prepare multi-scale information
\end{enumerate}

\subsection{The Different Constructors}

Frame has multiple constructors for different scenarios:

\begin{table}[H]
\centering
\begin{tabular}{@{}ll@{}}
\toprule
\textbf{Constructor} & \textbf{Use Case} \\
\midrule
\texttt{Frame()} & Default, empty frame \\
\texttt{Frame(const Frame\&)} & Copy constructor \\
\texttt{Frame(imLeft, imRight, ...)} & Stereo cameras \\
\texttt{Frame(imGray, imDepth, ...)} & RGB-D cameras \\
\texttt{Frame(imGray, timestamp, ...)} & Monocular camera \\
\texttt{Frame(imGray, keypoints, descriptors, ...)} & Pure XFeat mode \\
\bottomrule
\end{tabular}
\caption{Frame constructor variants}
\end{table}

\subsection{Monocular vs. XFeat Constructor}

The standard constructor calls \texttt{ExtractORB()}, which is slow.

Your custom \textbf{Pure XFeat constructor} skips this:

\begin{table}[H]
\centering
\begin{tabular}{@{}lcc@{}}
\toprule
\textbf{Step} & \textbf{Standard ORB} & \textbf{Pure XFeat} \\
\midrule
Feature Extraction & Calls \texttt{ExtractORB()} & Skipped (pre-extracted) \\
Undistortion & Calls \texttt{UndistortKeyPoints()} & Skipped (pre-undistorted) \\
Descriptors & 256-bit binary & 64-float vectors \\
Scale Levels & From ORBextractor (8 levels) & Manually set (3 levels) \\
\bottomrule
\end{tabular}
\caption{Comparison of Frame constructors}
\end{table}

\subsection{The Spatial Grid}

To speed up feature matching, the Frame divides the image into a grid:

\begin{lstlisting}[caption={Spatial grid for fast lookup}]
#define FRAME_GRID_COLS 64
#define FRAME_GRID_ROWS 48

// Each cell contains indices of keypoints in that region
vector<size_t> mGrid[FRAME_GRID_COLS][FRAME_GRID_ROWS];
\end{lstlisting}

This allows fast queries like: ``Find all features within 20 pixels of point (x, y)'' without scanning all N keypoints.

\subsection{Key Methods}

\begin{lstlisting}[caption={Important Frame methods}]
// Set/Get the camera pose
void SetPose(const Sophus::SE3f &Tcw);
Sophus::SE3f GetPose() const;

// Compute Bag-of-Words (for loop closure)
void ComputeBoW();

// Find features in a region (uses the grid)
vector<size_t> GetFeaturesInArea(float x, float y, float r);

// Check if a MapPoint is visible
bool isInFrustum(MapPoint* pMP, float viewingCosLimit);
\end{lstlisting}

\subsection{Frame vs KeyFrame}

\begin{table}[H]
\centering
\begin{tabular}{@{}lcc@{}}
\toprule
\textbf{Aspect} & \textbf{Frame} & \textbf{KeyFrame} \\
\midrule
Lifespan & Temporary & Permanent \\
Storage & Only current/last & All in map \\
Connections & None & Covisibility graph \\
Optimization & Pose only & Pose + structure \\
Loop closure & Cannot trigger & Can trigger \\
Count per sequence & Thousands & Hundreds \\
\bottomrule
\end{tabular}
\caption{Frames are temporary; KeyFrames are permanent}
\end{table}

%==============================================================================
\section{Act V: The Builder (\texttt{LocalMapping})}
%==============================================================================

While Tracking estimates poses, \texttt{LocalMapping} builds the map in the background.

\subsection{Why a Separate Thread?}

Local Bundle Adjustment (LBA) is too slow to run in real-time. If we did LBA for every frame, the system would lag behind. Instead:

\begin{itemize}
    \item \textbf{Tracking} runs fast (real-time) and just estimates pose
    \item \textbf{LocalMapping} runs slower but builds a better map
    \item They communicate through a queue of KeyFrames
\end{itemize}

\subsection{The Local Mapping Loop}

\begin{lstlisting}[caption={LocalMapping::Run() simplified}]
void LocalMapping::Run()
{
    while(!mbFinishRequested)
    {
        // Wait for new KeyFrame from Tracking
        if(CheckNewKeyFrames())
        {
            // 1. Process the new KeyFrame
            ProcessNewKeyFrame();
            
            // 2. Remove bad MapPoints
            MapPointCulling();
            
            // 3. Create new MapPoints by triangulation
            CreateNewMapPoints();
            
            // 4. Fuse duplicate MapPoints
            SearchInNeighbors();
            
            // 5. Local Bundle Adjustment
            if(!CheckNewKeyFrames())
            {
                Optimizer::LocalBundleAdjustment(mpCurrentKeyFrame, ...);
            }
            
            // 6. Remove redundant KeyFrames
            KeyFrameCulling();
        }
    }
}
\end{lstlisting}

\subsection{Step 1: ProcessNewKeyFrame}

When a new KeyFrame arrives:

\begin{lstlisting}[caption={ProcessNewKeyFrame}]
void LocalMapping::ProcessNewKeyFrame()
{
    // 1. Get the KeyFrame from the queue
    mpCurrentKeyFrame = mlNewKeyFrames.front();
    mlNewKeyFrames.pop_front();
    
    // 2. Compute Bag-of-Words
    mpCurrentKeyFrame->ComputeBoW();
    
    // 3. Associate MapPoints to the new KeyFrame
    for(int i = 0; i < mpCurrentKeyFrame->N; i++)
    {
        MapPoint* pMP = mpCurrentKeyFrame->mvpMapPoints[i];
        if(pMP)
            pMP->AddObservation(mpCurrentKeyFrame, i);
    }
    
    // 4. Update covisibility graph
    mpCurrentKeyFrame->UpdateConnections();
    
    // 5. Insert into the map
    mpAtlas->GetCurrentMap()->AddKeyFrame(mpCurrentKeyFrame);
}
\end{lstlisting}

\subsection{Step 3: CreateNewMapPoints}

This is where new 3D points are born:

\begin{lstlisting}[caption={CreateNewMapPoints simplified}]
void LocalMapping::CreateNewMapPoints()
{
    // For each neighboring KeyFrame
    for(KeyFrame* pKF2 : vpNeighKFs)
    {
        // Match features between current KF and neighbor
        vector<pair<int, int>> vMatchedIndices;
        ORBmatcher::SearchForTriangulation(mpCurrentKeyFrame, pKF2, 
                                           vMatchedIndices);
        
        // For each match, triangulate a new 3D point
        for(auto& match : vMatchedIndices)
        {
            cv::Point3f x3D = Triangulate(match);
            
            // Create new MapPoint
            MapPoint* pMP = new MapPoint(x3D, mpCurrentKeyFrame);
            pMP->AddObservation(mpCurrentKeyFrame, match.first);
            pMP->AddObservation(pKF2, match.second);
            
            mpAtlas->GetCurrentMap()->AddMapPoint(pMP);
        }
    }
}
\end{lstlisting}

\subsection{Step 5: Local Bundle Adjustment}

This is the expensive step that refines everything:

\begin{lstlisting}[caption={Local BA concept}]
void Optimizer::LocalBundleAdjustment(KeyFrame* pKF)
{
    // Collect local KeyFrames (connected to pKF)
    list<KeyFrame*> lLocalKeyFrames;
    
    // Collect local MapPoints (seen by local KFs)
    list<MapPoint*> lLocalMapPoints;
    
    // Collect fixed KeyFrames (neighbors of local KFs)
    list<KeyFrame*> lFixedKeyFrames;
    
    // Build g2o optimization graph
    g2o::SparseOptimizer optimizer;
    
    // Add pose vertices (local KFs)
    // Add point vertices (local MPs)
    // Add edges (reprojection errors)
    
    // Optimize!
    optimizer.optimize(5);  // 5 iterations
    
    // Update KeyFrame poses and MapPoint positions
}
\end{lstlisting}

\subsection{Step 6: KeyFrame Culling}

Remove redundant KeyFrames to keep the map compact:

\begin{lstlisting}[caption={KeyFrameCulling concept}]
void LocalMapping::KeyFrameCulling()
{
    // For each local KeyFrame
    for(KeyFrame* pKF : mvpLocalKeyFrames)
    {
        // Count how many of its MapPoints are also seen
        // by at least 3 other KeyFrames
        int nRedundantObs = 0;
        
        for(MapPoint* pMP : pKF->mvpMapPoints)
        {
            if(pMP->Observations() >= 3)
                nRedundantObs++;
        }
        
        // If 90% of points are redundant, cull the KeyFrame
        if(nRedundantObs > 0.9 * pKF->TrackedMapPoints())
        {
            pKF->SetBadFlag();  // Mark for deletion
        }
    }
}
\end{lstlisting}

\subsection{Communication with Tracking}

The interaction is through shared state:

\begin{lstlisting}[caption={Tracking to LocalMapping communication}]
// In Tracking::CreateNewKeyFrame()
mpLocalMapper->InsertKeyFrame(pKF);

// In LocalMapping
void InsertKeyFrame(KeyFrame* pKF)
{
    unique_lock<mutex> lock(mMutexNewKFs);
    mlNewKeyFrames.push_back(pKF);
}
\end{lstlisting}

%==============================================================================
\section{Act VI: The Proofreader (\texttt{LoopClosing})}
%==============================================================================

As the camera moves, small errors accumulate (drift). \texttt{LoopClosing} detects when the camera returns to a previously visited place and corrects the drift.

\subsection{Why Loop Closure Matters}

Imagine walking around a building and returning to your starting point:
\begin{itemize}
    \item Without loop closure: The map shows you're 10 meters away from the start
    \item With loop closure: The system recognizes you're back at start and fixes the error
\end{itemize}

\subsection{The Loop Closing Pipeline}

\begin{lstlisting}[caption={LoopClosing::Run() simplified}]
void LoopClosing::Run()
{
    while(!mbFinishRequested)
    {
        if(CheckNewKeyFrames())
        {
            // 1. Detect loop candidates using Bag-of-Words
            if(DetectLoop())
            {
                // 2. Compute similarity transformation (with scale)
                if(ComputeSim3())
                {
                    // 3. Correct the loop
                    CorrectLoop();
                    
                    // 4. Global Bundle Adjustment (in a new thread)
                    RunGlobalBundleAdjustment();
                }
            }
        }
    }
}
\end{lstlisting}

\subsection{Step 1: Loop Detection}

We use Bag-of-Words to find similar KeyFrames:

\begin{lstlisting}[caption={DetectLoop concept}]
bool LoopClosing::DetectLoop()
{
    // 1. Compute minimum score threshold
    //    (use similarity with neighbors as baseline)
    float minScore = ComputeMinScoreInNeighbors();
    
    // 2. Query database for similar KeyFrames
    vector<KeyFrame*> vpCandidates = 
        mpKeyFrameDB->DetectLoopCandidates(mpCurrentKF, minScore);
    
    if(vpCandidates.empty())
        return false;
    
    // 3. Check temporal consistency
    //    (candidates should appear in multiple consecutive KFs)
    for(KeyFrame* pKF : vpCandidates)
    {
        if(CheckConsistency(pKF))
            mvpEnoughConsistentCandidates.push_back(pKF);
    }
    
    return !mvpEnoughConsistentCandidates.empty();
}
\end{lstlisting}

\subsection{Step 2: Compute Sim3}

We need to find the transformation between current and loop KeyFrame:

\begin{lstlisting}[caption={ComputeSim3 concept}]
bool LoopClosing::ComputeSim3()
{
    for(KeyFrame* pKF : mvpEnoughConsistentCandidates)
    {
        // 1. Match features between current KF and candidate
        vector<MapPoint*> vpMatchedPoints;
        int nMatches = ORBmatcher::SearchByBoW(mpCurrentKF, pKF, 
                                               vpMatchedPoints);
        
        if(nMatches < 20)
            continue;
        
        // 2. Compute Sim3 using RANSAC
        Sim3Solver solver(mpCurrentKF, pKF, vpMatchedPoints);
        
        cv::Mat Scm;  // Similarity transformation
        if(solver.ComputeSim3(Scm))
        {
            mScw = Scm;
            mpMatchedKF = pKF;
            return true;
        }
    }
    return false;
}
\end{lstlisting}

\textbf{Note:} \texttt{Sim3} includes scale, which is important for monocular SLAM where scale can drift.

\subsection{Step 3: Correct Loop}

Once we find a valid loop, we correct the drift:

\begin{lstlisting}[caption={CorrectLoop concept}]
void LoopClosing::CorrectLoop()
{
    // 1. Stop LocalMapping temporarily
    mpLocalMapper->RequestStop();
    
    // 2. Compute correction for current KeyFrame
    Sophus::Sim3f Scw = mpCurrentKF->GetPose() * mScw.inverse();
    
    // 3. Propagate correction to all connected KeyFrames
    for(KeyFrame* pKF : mvpCurrentConnectedKFs)
    {
        Sophus::Sim3f CorrectedPose = Scw * pKF->GetPoseRelativeTo(mpCurrentKF);
        pKF->SetPose(CorrectedPose);
    }
    
    // 4. Correct all associated MapPoints
    for(MapPoint* pMP : mvpLoopMapPoints)
    {
        pMP->SetWorldPos(CorrectedPos);
    }
    
    // 5. Fuse duplicate MapPoints
    SearchAndFuse();
    
    // 6. Pose-graph optimization
    Optimizer::OptimizeEssentialGraph(...);
    
    // 7. Resume LocalMapping
    mpLocalMapper->Release();
}
\end{lstlisting}

\subsection{Step 4: Global Bundle Adjustment}

After loop closure, run global BA to refine everything:

\begin{lstlisting}[caption={Global BA}]
void LoopClosing::RunGlobalBundleAdjustment()
{
    // Run in a separate thread to not block loop closing
    thread* pThread = new thread(&LoopClosing::GlobalBundleAdjustmentThread, 
                                 this);
}

void LoopClosing::GlobalBundleAdjustmentThread()
{
    // Optimize ALL KeyFrames and MapPoints
    Optimizer::GlobalBundleAdjustment(mpAtlas->GetCurrentMap());
    
    // Merge corrected map back
    // ...
}
\end{lstlisting}

\subsection{The Covisibility Graph}

Loop closing uses the covisibility graph to propagate corrections:

\begin{itemize}
    \item \textbf{Nodes}: KeyFrames
    \item \textbf{Edges}: Weighted by number of shared MapPoints
    \item \textbf{Use}: Find which KeyFrames are affected by a correction
\end{itemize}

\begin{lstlisting}[caption={Covisibility example}]
// Get KeyFrames that share at least 15 MapPoints with pKF
vector<KeyFrame*> vpCovisibleKFs = pKF->GetCovisiblesByWeight(15);
\end{lstlisting}

%==============================================================================
\section{Epilogue: The Complete Picture}
%==============================================================================

\subsection{Data Flow Summary}

\begin{verbatim}
                 IMAGE
                   |
                   v
    +---------------------------+
    |      mono_euroc.cc        |  Load image, call TrackMonocular
    +---------------------------+
                   |
                   v
    +---------------------------+
    |         System            |  Convert to grayscale, delegate
    +---------------------------+
                   |
                   v
    +---------------------------+
    |        Tracking           |  Create Frame, estimate pose
    +---------------------------+
          |              \
          v               v
    +-----------+   +---------------+
    |   Frame   |   | LocalMapping  |  (Background)
    +-----------+   +---------------+
          |               |
          v               v
    +-----------+   +---------------+
    | MapPoint  |   | LoopClosing   |  (Background)
    +-----------+   +---------------+
\end{verbatim}

\subsection{Thread Model}

\begin{table}[H]
\centering
\begin{tabular}{@{}lll@{}}
\toprule
\textbf{Thread} & \textbf{Speed} & \textbf{Responsibility} \\
\midrule
Main (Tracking) & Real-time (30 fps) & Pose estimation \\
LocalMapping & Background & Map building, local BA \\
LoopClosing & Background & Loop detection, global BA \\
Viewer & Background & Visualization \\
\bottomrule
\end{tabular}
\caption{ORB-SLAM3 threading model}
\end{table}

\subsection{The Four Core Data Structures}

\begin{enumerate}
    \item \textbf{Frame}: A snapshot with features (temporary)
    \item \textbf{KeyFrame}: An important snapshot stored permanently
    \item \textbf{MapPoint}: A 3D point in the world
    \item \textbf{Map/Atlas}: Collection of KeyFrames and MapPoints
\end{enumerate}

\subsection{Where to Go Next}

Now that you understand the big picture, you can dive deeper into:

\begin{enumerate}
    \item \textbf{Feature Matching}: How does \texttt{ORBmatcher} work?
    \item \textbf{Optimization}: How does g2o refine the poses?
    \item \textbf{Map Management}: How does Atlas handle multiple maps?
    \item \textbf{IMU Integration}: How are inertial measurements fused?
    \item \textbf{Your XFeat Modifications}: How does the Pure XFeat pipeline differ?
\end{enumerate}

\subsection{Key Takeaways}

\begin{itemize}
    \item ORB-SLAM3 is a \textbf{multi-threaded} system with clear separation of concerns
    \item \textbf{Tracking} is the hot path that must run in real-time
    \item \textbf{LocalMapping} and \textbf{LoopClosing} do the heavy lifting in the background
    \item Everything revolves around \textbf{KeyFrames} and \textbf{MapPoints}
    \item The \textbf{covisibility graph} connects KeyFrames that see the same points
    \item \textbf{Bundle Adjustment} is the secret sauce that makes it all accurate
\end{itemize}


% Appendices
\appendix
%==============================================================================
\section{Appendix A: File Organization}
%==============================================================================

\subsection{Project Structure}

\begin{verbatim}
ORB_SLAM3/
|
+-- Examples/                    <-- Entry points
|   +-- Monocular/
|   |   +-- mono_euroc.cc        <-- Main for EuRoC dataset
|   |   +-- mono_tum.cc          <-- Main for TUM dataset
|   +-- Stereo/
|   +-- RGB-D/
|
+-- include/                     <-- Header files
|   +-- System.h                 <-- The orchestrator
|   +-- Tracking.h               <-- Pose estimation
|   +-- LocalMapping.h           <-- Map building
|   +-- LoopClosing.h            <-- Drift correction
|   +-- Frame.h                  <-- Single image
|   +-- KeyFrame.h               <-- Important image
|   +-- MapPoint.h               <-- 3D point
|   +-- Map.h                    <-- Holds KFs and MPs
|   +-- Atlas.h                  <-- Multi-map management
|   +-- ORBextractor.h           <-- Feature extraction
|   +-- ORBmatcher.h             <-- Feature matching
|   +-- Optimizer.h              <-- g2o wrapper
|   +-- GeometricCamera.h        <-- Camera models
|
+-- src/                         <-- Implementation files
|   +-- (mirrors include/)
|
+-- Thirdparty/                  <-- External libraries
|   +-- DBoW2/                   <-- Bag-of-Words
|   +-- g2o/                     <-- Graph optimization
|   +-- Sophus/                  <-- Lie groups (SE3)
|
+-- Vocabulary/
|   +-- ORBvoc.txt               <-- Visual vocabulary (~100MB)
|
+-- docs/
    +-- Understanding_code/      <-- This document!
\end{verbatim}

\subsection{Your Custom Files}

\begin{verbatim}
ORB_SLAM3_perso_Branch_ORB_LOOP/
|
+-- include/
|   +-- ONNXFeatureExtractor.h   <-- XFeat extraction
|   +-- LightGlueMatcher.h       <-- LightGlue matching
|   +-- ONNXDepthEstimator.h     <-- Depth estimation
|   +-- SemiDenseTracker.h       <-- Semi-dense tracking
|
+-- src/
    +-- ONNXFeatureExtractor.cc
    +-- LightGlueMatcher.cc
    +-- ...
\end{verbatim}

\subsection{Key Configuration Files}

\begin{verbatim}
Examples/Monocular/
+-- EuRoC.yaml                   <-- Standard EuRoC config
+-- TUM_mono.yaml                <-- TUM config
+-- TUM_PureXFeat.yaml           <-- Your XFeat config
\end{verbatim}

\subsection{Build Artifacts}

\begin{verbatim}
build/                           <-- CMake build directory
+-- lib/
|   +-- libORB_SLAM3.so          <-- Shared library
+-- Examples/Monocular/
    +-- mono_euroc               <-- Executable
\end{verbatim}

%==============================================================================
\section{Appendix B: Quick Reference}
%==============================================================================

\subsection{Key Data Types}

\begin{table}[H]
\centering
\begin{tabular}{@{}llp{7cm}@{}}
\toprule
\textbf{Type} & \textbf{Variable} & \textbf{Meaning} \\
\midrule
\texttt{Sophus::SE3f} & \texttt{mTcw} & Camera pose (world $\to$ camera) \\
\texttt{Sophus::SE3f} & \texttt{mTwc} & Camera pose (camera $\to$ world) \\
\texttt{Eigen::Vector3f} & \texttt{mOw} & Camera center in world coords \\
\texttt{Eigen::Matrix3f} & \texttt{mRcw} & Rotation (world $\to$ camera) \\
\texttt{cv::Mat} & \texttt{mK} & Camera calibration matrix ($3\times3$) \\
\texttt{cv::Mat} & \texttt{mDistCoef} & Distortion coefficients \\
\texttt{cv::Mat} & \texttt{mDescriptors} & Feature descriptors ($N\times32$ for ORB) \\
\texttt{vector<cv::KeyPoint>} & \texttt{mvKeys} & 2D keypoint locations \\
\texttt{vector<cv::KeyPoint>} & \texttt{mvKeysUn} & Undistorted keypoints \\
\texttt{vector<MapPoint*>} & \texttt{mvpMapPoints} & 3D associations for each keypoint \\
\texttt{DBoW2::BowVector} & \texttt{mBowVec} & Bag-of-Words representation \\
\bottomrule
\end{tabular}
\caption{Common variables you'll encounter}
\end{table}

\subsection{Naming Conventions}

\begin{table}[H]
\centering
\begin{tabular}{@{}ll@{}}
\toprule
\textbf{Prefix} & \textbf{Meaning} \\
\midrule
\texttt{m} & Member variable \\
\texttt{mp} & Member pointer \\
\texttt{mv} & Member vector \\
\texttt{mb} & Member boolean \\
\texttt{mn} & Member number (int/unsigned) \\
\texttt{mf} & Member float \\
\texttt{p} & Pointer (local variable) \\
\texttt{v} & Vector (local variable) \\
\texttt{b} & Boolean (local variable) \\
\texttt{n} & Number (local variable) \\
\bottomrule
\end{tabular}
\caption{Hungarian notation used in ORB-SLAM3}
\end{table}

\subsection{Key Functions Cheat Sheet}

\subsubsection{System}
\begin{itemize}
    \item \texttt{TrackMonocular(im, timestamp)} $\rightarrow$ Process one image
    \item \texttt{Shutdown()} $\rightarrow$ Stop all threads
    \item \texttt{SaveTrajectoryEuRoC(file)} $\rightarrow$ Save poses
\end{itemize}

\subsubsection{Tracking}
\begin{itemize}
    \item \texttt{GrabImageMonocular(im, t)} $\rightarrow$ Entry point
    \item \texttt{Track()} $\rightarrow$ Main logic
    \item \texttt{TrackWithMotionModel()} $\rightarrow$ Fast pose estimation
    \item \texttt{TrackReferenceKeyFrame()} $\rightarrow$ Fallback
    \item \texttt{TrackLocalMap()} $\rightarrow$ Refine pose
    \item \texttt{CreateNewKeyFrame()} $\rightarrow$ Insert KF
\end{itemize}

\subsubsection{Frame}
\begin{itemize}
    \item \texttt{SetPose(Tcw)} $\rightarrow$ Set camera pose
    \item \texttt{GetPose()} $\rightarrow$ Get camera pose
    \item \texttt{ComputeBoW()} $\rightarrow$ Compute Bag-of-Words
    \item \texttt{GetFeaturesInArea(x, y, r)} $\rightarrow$ Spatial query
\end{itemize}

\subsubsection{Optimizer}
\begin{itemize}
    \item \texttt{PoseOptimization(\&frame)} $\rightarrow$ Refine frame pose
    \item \texttt{LocalBundleAdjustment(pKF)} $\rightarrow$ Local BA
    \item \texttt{GlobalBundleAdjustment(pMap)} $\rightarrow$ Global BA
\end{itemize}

\subsection{Common Debug Patterns}

\begin{lstlisting}[caption={Useful debug prints}]
// Print current frame ID and timestamp
cout << "Frame " << mCurrentFrame.mnId 
     << " @ " << mCurrentFrame.mTimeStamp << endl;

// Print number of features
cout << "Features: " << mCurrentFrame.N << endl;

// Print number of matched MapPoints
int nMatches = 0;
for(int i = 0; i < mCurrentFrame.N; i++)
    if(mCurrentFrame.mvpMapPoints[i])
        nMatches++;
cout << "Matches: " << nMatches << endl;

// Print camera pose
cout << "Pose:\n" << mCurrentFrame.GetPose().matrix() << endl;
\end{lstlisting}

\subsection{Key Thresholds}

\begin{table}[H]
\centering
\begin{tabular}{@{}lll@{}}
\toprule
\textbf{Parameter} & \textbf{Typical Value} & \textbf{Used In} \\
\midrule
Min matches for init & 100 & MonocularInitialization \\
Min matches for tracking & 20 & TrackWithMotionModel \\
Min inliers for local map & 30 & TrackLocalMap \\
BoW similarity threshold & 0.8 & Loop detection \\
Triangulation parallax & 1.0° & CreateNewMapPoints \\
\bottomrule
\end{tabular}
\caption{Important thresholds in ORB-SLAM3}
\end{table}


\end{document}
